\section{Conclusion}

The paper now establishes both sides of the frontier question. On the positive side, the current formalized tractable landscape factors into primitive mechanisms, lifts, and degeneracies. On the negative side, quotient realizability is maximal, so quotient shape cannot classify tractability by itself, and an admissibility-layer collapse theorem rules out naive finite classifiers even after imposing theorem-forced closure laws together with computational and structural restrictions.

Technically, the decisive mechanism is closure-orbit disagreement via pair-targeted action-independent affine transport. The same template now applies to four obstruction families (dominant pair, margin masking, ghost action, and additive/statewise offset), yielding a formal admissible-collapse theorem.

The open problem is therefore sharper than before. The question is no longer whether simple finite lists or unconstrained quotient predicates can characterize the boundary; they cannot. The research program is redirected toward a narrower target: identify a stronger representation-sensitive principle, outside the present admissible surface but still mathematically natural, that could support a genuine optimizer-compatible dichotomy. In that sense, the meta-impossibility theorem is not only a negative result. It determines which kinds of frontier theorem are no longer worth pursuing.
